\chapter{Machines, Calling Conventions, and Object Files}\label{ch:machines-conventions-object-files}

The middle end works with an idealized machine that has unlimited temporaries and no fixed call rules.
The back end has no such luxury.
It must target real registers, real stacks, real argument-passing rules, and real file formats that loaders understand.
This chapter builds that bridge.
It explains what a target machine looks like to a compiler, how callers and callees agree on data exchange, and how object files carry code and metadata to the linker.

\section{From IR to a Real Machine}\label{sec:ch24-machine-model}

A compiler models a processor as a small set of facts: the register file, the instruction set, the addressing modes, the stack discipline, and the assembler and object format it must emit.
\term{Register} means a fast storage location named by instructions, grouped into classes such as integer, floating-point, and vector.
\term{Calling convention} means the contract that assigns arguments and return values to registers or stack slots, designates callee-saved versus caller-saved registers, and fixes stack alignment and cleanup duties.
\term{Application binary interface} (ABI) extends that contract to type layout, structure padding, exception tables, and object-file details so separately compiled units interoperate.
Two compilers that agree on the source language but disagree on the ABI produce objects that crash when linked together.

Consider adding two 64-bit integers on a load-store machine with sixteen general-purpose registers.
The IR says \code{t3 = add t1, t2} with no commitment about locations.
The machine needs each operand in a register, the result in a register, and any spilled value reachable through a stack slot.
If the values arrive as function arguments, the calling convention decides whether they start in registers or on the stack.
The rest of the back end, selection, scheduling, and allocation, exists to close exactly this gap.
A \term{stack frame} holds the return address, saved registers, spilled locals, outgoing argument space, and alignment padding for one activation.
A \term{red zone} is a small area below the stack pointer that a leaf function may use without adjusting the pointer, where the ABI provides one.

The frame layout below makes the abstract concrete for a typical 64-bit target.
\begin{irlisting}
; abstract frame for f(a, b) with one spill
frame f:
  incoming a -> reg r0, incoming b -> reg r1
  spill s1 -> [sp + 0], alignment pad -> [sp + 8]
  save ra -> [sp + 16], save fp -> [sp + 24]
  call g: move args to r0, r1; call g; result in r0
\end{irlisting}
Every line of that sketch corresponds to a decision the back end must make explicit.
Which registers carry \code{a} and \code{b} on entry.
Where the spill lives and how it is addressed.
Who saves the return address and frame pointer.
Who cleans the outgoing argument space after the call to \code{g}.
The calling convention answers each question the same way for every function so that separate translation units agree.

\section{Calling Conventions in Detail}\label{sec:ch24-conventions}

Argument passing follows a fixed algorithm: classify each argument, assign registers while they last, and spill the rest to the stack in order.
Integer and pointer arguments usually take one register class, floating-point arguments take another, and large structures are passed by hidden pointer, by value copy, or split across registers depending on the ABI.
Return values use a dedicated register or register pair, with a hidden first parameter for large returns.
\term{Caller-saved register} means the caller preserves the value across a call if it still needs it.
\term{Callee-saved register} means the callee preserves it on behalf of all callers.
\term{Stack alignment} forces the stack pointer to a boundary such as 16 bytes at each call so vector instructions and system calls behave correctly.
Figure~\ref{fig:ch24-frame} shows a canonical frame with each region in place.
Table~\ref{tab:ch24-convention-choices} compares the main convention decisions and their costs.

\begin{figure}[htbp]
\centering
\begin{tikzpicture}[node distance=0.55cm]
\node[basic block] (high) {incoming args, saved fp/ra};
\node[basic block, below=of high] (mid) {callee-saved regs};
\node[basic block, below=of mid] (spill) {spills and locals};
\node[basic block, below=of spill] (out) {outgoing arg space};
\draw[cfg edge] (high) -- (mid);
\draw[cfg edge] (mid) -- (spill);
\draw[cfg edge] (spill) -- (out);
\end{tikzpicture}
\caption{A canonical stack frame grows from saved state through spills to outgoing argument space, with alignment padding inserted where the ABI requires it.}
\label{fig:ch24-frame}
\end{figure}

\begin{table}[htbp]
\caption{Calling-convention decisions and what each choice costs.}
\label{tab:ch24-convention-choices}
\begin{tabular}{lll}
\toprule
Decision & Register-heavy choice & Stack-heavy choice \\
\midrule
Argument passing & More registers, less memory traffic & Simpler ABI, more loads \\
Callee-saved set & Fewer caller spills & Larger prologues \\
Return by value & Registers for small types & Hidden pointer, extra copy \\
Variadic calls & Fixed register sequence & Stack overflow area \\
Alignment & 16-byte boundary always & Smaller but fragile frames \\
\bottomrule
\end{tabular}
\end{table}

Variadic functions deserve attention because the callee does not know its arity statically.
The convention therefore fixes how the register arguments are also mirrored to the stack or counted in a register so that \code{va\_start} can find them.
Floating-point and vector rules interact with this design: some ABIs pass a count of vector registers used, others require the caller to store register arguments home.
Structure passing is the other classic trap.
A two-word structure may split across two registers on one ABI and go entirely on the stack on another.
Bit-fields, homogeneous aggregates, and complex numbers each have their own classification pages in real ABI documents.
The compiler implements the classification exactly, because one wrong byte offset corrupts every cross-module call.
\begin{ccnote}[Read the ABI as a specification]
ABI documents read like standards because they are standards.
Implement the classification algorithm case by case and test it against the platform conformance suite.
Guessing from one example struct invites silent miscompilation.
\end{ccnote}

Caller-saved and callee-saved choices trade prologue cost against call-site cost.
A leaf function that uses few registers prefers caller-saved temporaries so it saves nothing.
A deeply nested non-leaf function prefers callee-saved registers whose save cost is paid once in its prologue rather than at every call it makes.
Register allocators know this split and bias their choices accordingly.
The stack unwinder knows it too: exception tables record which registers were saved and where, so the runtime can restore them during propagation.
That is why frame layout, debug information, and unwinding tables are generated together rather than by three independent passes.

\section{Lowering Calls and Frames}\label{sec:ch24-algorithm}

Algorithm~\ref{alg:ch24-lower-call} turns a machine-independent call into explicit moves, saves, and stack adjustments.
It runs during instruction selection or just after, when the compiler still sees virtual registers but already knows the target ABI.
The input is a call node with a typed argument list and a function signature.
The output is a lowered sequence with a resolved frame layout for the enclosing function.
Preconditions matter: types must already be laid out, structures classified, and the callee type visible enough to pick the right convention variant.

\begin{algorithm}[htbp]
\caption{Lower a call and fix the enclosing frame}\label{alg:ch24-lower-call}
\begin{pseudocode}
\Require call node $C$ with signature $S$, current frame $F$
\Ensure ABI-conformant sequence replacing $C$, updated $F$
\State classify each argument of $S$ into registers or stack slots
\State assign return location: register pair or hidden pointer
\State spill all live caller-saved registers around $C$
\State reserve the outgoing area and establish the ABI stack position and alignment
\State store stack arguments at their ABI offsets in $F$
\State place register arguments with a parallel copy; break copy cycles using a scratch location
\State verify stack alignment at the call point
\State emit the call and copy the result to a fresh virtual
\State restore spilled caller-saved registers
\State record unwind and debug info for new frame ranges
\end{pseudocode}
\end{algorithm}

The algorithm is linear in the number of arguments plus the number of live caller-saved values, hence $O(A + L)$ per call site, where $A$ is the argument count and $L$ is the live set size at the site.
Whole-function frame assignment costs $O(N + E)$ for $N$ instructions and $E$ dataflow edges when liveness is available.
The expensive part is not the classification but the spill and restore traffic it implies, which the register allocator later tries to reduce.
The pseudocode combines logical lowering with eventual physical-register repair for exposition.
Before allocation, represent argument locations, return locations, and call clobbers as constraints; after allocation, save only values live across the call whose assigned registers it may destroy.
Move the result to a safe location before restoring a saved register that overlaps the return channel.
A common refinement coalesces the argument moves with the producer instructions so that values are born in their ABI registers.
Another refinement shares outgoing argument space among call sites instead of sizing it for the worst site plus padding each time.

Several edge cases need explicit handling.
Tail calls reuse the current frame when the callee signature permits, which the lowering pass must verify rather than assume.
Indirect calls load the target into a designated register and may need an extra guard page or branch-target annotation on hardened platforms.
Calls to dynamically resolved symbols route through a procedure-linkage stub instead of a direct PC-relative jump.
Each of these variants changes the instruction pattern but not the skeleton of the algorithm.
% TODO(cite): need a primary reference for modern ABI frame layout rules.

\section{Worked Example}\label{sec:ch24-worked}

We lower a concrete call by hand for a RISC-like target with eight argument registers \code{a0} through \code{a7}, return register \code{rv}, and 16-byte stack alignment.
The source call is \code{r = f(10, 20, 30, 40, 50, 60, 70, 80, 90)} with nine integer arguments.
The first eight fit in registers and the ninth goes on the stack.
Assume no caller-saved values are live, so no spill code appears.
\begin{irlisting}
; before lowering
r = call f(10, 20, 30, 40, 50, 60, 70, 80, 90)
; after lowering
mov a0, 10
mov a1, 20
mov a2, 30
mov a3, 40
mov a4, 50
mov a5, 60
mov a6, 70
mov a7, 80
store [sp + 0], 90
call f
mov r, rv
\end{irlisting}
Now check alignment.
Suppose the stack pointer is 16-byte aligned on entry and the frame already reserves 8 bytes for a spill.
Storing one 8-byte argument makes the used space 16 bytes, so alignment holds.
If the frame had reserved 0 bytes instead, the lowering pass would insert 8 bytes of padding before the store.
That padding is invisible in the IR but mandatory on the machine.

Next add a callee-saved wrinkle.
Suppose \code{f} itself uses two callee-saved registers \code{s0} and \code{s1} plus one 8-byte local.
Its prologue stores \code{s0} and \code{s1} (16 bytes) and reserves 8 bytes for the local plus 8 bytes of padding to keep 16-byte alignment, for a 32-byte frame.
Its epilogue restores the two registers and releases the frame before returning.
The caller never saves \code{s0} or \code{s1}; the contract guarantees their values survive the call.
If the allocator had picked caller-saved registers instead, each live value would need an explicit spill and reload around the call.
Counting those bytes by hand shows why register-class choice is a size and speed decision, not just a naming decision.

\section{Implementation Notes}\label{sec:ch24-implementation}

Production back ends isolate machine knowledge in a target description rather than scattering conditionals through the code.
LLVM uses TableGen records for registers, classes, calling conventions, and instruction patterns, while GCC uses machine-description files with a similar role.
Both generate the classifier, the frame-layout logic, and parts of the selector from those descriptions.
Adding a new ABI variant then means editing data plus a few hooks, not rewriting every call site handler.
The complexity of the description is real, but the alternative of hand-coded per-function logic is worse.

Object files close the loop between compiler and linker.
An \term{object file} holds machine code sections, data sections, a symbol table, and relocations that name the fixups the linker or loader must apply.
A \term{relocation} records an offset, a symbol, and a computation such as \enquote{add the symbol address plus addend} so the linker can patch the bytes once addresses are known.
Common formats include ELF on Unix-like systems, Mach-O on Apple platforms, and COFF on Windows.
Each carries the same ideas under different section names plus ABI tags, unwind tables, and debug sections.
The compiler must emit the right directives for symbol visibility, thread-local storage, position independence, and exception tables, or linking succeeds and the program still fails at load time.

Testing this layer means checking bytes and boundaries, not just program output.
Conformance tests pass structures by value, exercise variadic calls, alternate register and stack arguments, and verify alignment at every call.
Unwind tests throw through inlined and optimized frames.
Link tests combine objects built with different optimization levels to catch ABI drift.
Useful background on machine-level code generation structure appears in~\cite{appel-1998-modern} and see~\cite{morgan-1998-optimizing}.
Correctness arguments for preserving observable behavior across lowering appear in~\cite{safonov-2010-trustworthy}.
The source-level type and layout rules that feed classification are defined in~\cite{wg14-2011-c11}.

\section{Exercises}\label{sec:ch24-exercises}

\begin{enumerate}
\item Lower the call \code{r = g(x, y, z)} on a target with two integer argument registers. \label{ex:ch24-lower-three} Assume all three arguments are live integers and no caller-saved values are live. Show the moves, the stack store, and the alignment check.
\item A function uses three callee-saved registers and has 24 bytes of locals. \label{ex:ch24-frame-size} Compute the smallest 16-byte-aligned frame size, showing padding, save area, and locals separately.
\item An object file references an external function \code{h}. \label{ex:ch24-reloc} List the symbol-table entry and the relocation the compiler must emit for a direct call, and explain what the linker patches.
\end{enumerate}

\section{Further Reading}\label{sec:ch24-reading}

Machine-level organization, frame management, and code-generator structure are covered in~\cite{appel-1998-modern}.
Calling-sequence design and target-specific optimization tradeoffs are treated in~\cite{morgan-1998-optimizing}.
Validation of lowering against observable behavior is discussed in~\cite{safonov-2010-trustworthy}.
Type layout, alignment, and argument-passing constraints at the language level appear in~\cite{wg14-2011-c11}.
